\documentclass[10pt,a4paper]{amsart}
\usepackage[margin=19mm]{geometry}
\usepackage{amsmath,amssymb,mathtools}
\usepackage[scr=rsfs,cal=euler]{mathalfa}
\usepackage{xfrac}
\usepackage[unicode,hidelinks,bookmarks=false]{hyperref}
\newcommand{\ie}{i.e.\ }
\newcommand{\cf}{cf.\ }
\newcommand{\resp}{respectively\ }
\newcommand{\Subsection}{\subsection}
\providecommand{\nobreakdash}{\nobreak\hbox{-}}
\setlength{\emergencystretch}{2em}
\multlinegap0pt
\usepackage{amssymb}
\usepackage[shortlabels]{enumitem}
\newtheorem{lemma}{Lemma}
\newcommand{\Q}{\mathbb{Q}}
\title{The classical topological invariants of homogeneous spaces}
\author{John Jones, Dmitriy Rumynin, Adam R. Thomas}
\date{Source: arXiv:2310.14365v2 (CC BY 4.0)}
\begin{document}
\maketitle

\noindent\emph{Source and licence.} The classical topological invariants of homogeneous spaces. Version arXiv:2310.14365v2 (CC BY 4.0), distributed by arXiv under the Creative Commons Attribution 4.0 International licence, http://creativecommons.org/licenses/by/4.0/, as recorded in the arXiv licence metadata for this exact version and shown on the abstract page https://arxiv.org/abs/2310.14365v2. Full licence notice: \texttt{licenses/arxiv-2310.14365-CC-BY-4.0.txt}.\par\smallskip

\noindent\textit{Editorial scope.} Counted unit: the article's lemma lem:cyclicpsimodules (source0.tex lines 759--764). The article's complete original proof of it is delivered with it (source0.tex lines 765--786, one \texttt{proof} environment of 22 lines). No mathematics is written, repaired or re-derived: the statement and the proof are the article's own lines, in the article's own notation, and no retained line is substituted or re-flowed.  No further source-local context is needed: every cross-reference of every retained line resolves inside the two delivered spans.  Nothing was omitted from any retained span, so the package carries no omission record.  No retained line contains a citation, so the wrapper carries no bibliography and no bibliography entry.  No retained line contains a figure environment, an image-inclusion command or drawing code, so no image is delivered and the package declares no asset dependency.  Editorial formatting only, in addition: an AMS \texttt{amsart} preamble that restates the article's own declarations and macros used by the retained lines, namely the environments \texttt{lemma} on separate counters, together with the standard packages those lines need.  The retained environments print in this document as Lemma 1 in order of appearance, whereas the article prints the same items with the numbering of its own full document.  The complete native source of the article as distributed in the arXiv e-print for this exact version is bundled unchanged as \texttt{sources/148885/source0.tex}.
\begin{lemma} \label{lem:cyclicpsimodules}
\begin{enumerate}
\item The $\Psi$-module $\Q(r_1, \dots, r_l)$ is a cyclic $\Psi$-module if and only if the $r_i$ are distinct.
\item Suppose $r_1, \dots, r_l$ are distinct. Then $u$ is a generator of the $\Psi$-module $\Q(r_1, \dots, r_l)$ if and only if $u = u_1 + \dots + u_l$ where $u_i$ is a non-zero element of $\Q(r_1, \dots, r_l)$ with weight $r_i$.
\end{enumerate}
\end{lemma}

\begin{proof} 
Let $r$ be one of the $r_i$ and let $W_r$ be the weight $r$ subspace of $\Q(r_1, \dots, r_l)$.  So the dimension of $W_r$ is the number of $r_i$ equal to $r$.  Let $p_r : \Q(r_1, \dots, r_l) \to W_r$ be the standard projection onto $W_r$.  Suppose $\Q(r_1, \dots, r_l)$ is a cyclic $\Psi$-module. Then, since $p_r$ is a map of $\Psi$-modules $W_r$ is also a cyclic $\Psi$-module.  Let $w$ be a non-zero element of $W_r$.  Then $\psi^kw = k^rw$ and so the $\Psi$-module generated by $w$ must be $1$-dimensional and the dimension of $W_r$ must be one.  It follows that the $r_1, \dots, r_l$ must be distinct.  



Recall that for a linear map $A : V \to V$, a vector $v \in V$ is called a cyclic vector for $A$ if 
$$
v, A(v), \dots, A^{l-1}(v)
$$
is a basis for $V$. 
 
The eigenvalues of $\psi^k$ are $k^{r_1}, \dots, k^{r_l}$.  If $r_1, \dots, r_l$ are distinct
and $k\geq 2$, 
then clearly %%it is straightforward to check that
there exists a cyclic vector for $\psi^k$. If $u$ is such a cyclic vector, it follows that
$$
u, \ \psi^k(u), \ \psi^{k^2}(u), \ldots, \ \psi^{k^{l-1}}(u)
$$ 
are a basis for $\Q(r_1, \dots, r_l)$ and so $\Q(r_1, \dots, r_l)$ is generated as a $\Psi$-module by $u$.

This completes the proof of the first part of the theorem. The second part follows directly from the above characterisation of cyclic vectors.
\end{proof}

\end{document}
